\documentclass[a4paper,12pt]{elsarticle}
\usepackage{amsmath}
\usepackage{algorithm}
\usepackage{ulem}
\usepackage{amssymb}
\usepackage[dvipsnames]{xcolor}
\usepackage[mathlines,displaymath]{lineno}
\usepackage{graphicx}
\usepackage{soul}
\usepackage[all,2cell,dvips]{xy}
\usepackage{natbib}
\usepackage[usenames,rgb]{xcolor} 
\usepackage{ucs}                                            
\usepackage[utf8]{inputenc}                                         
\usepackage{minted}
\usepackage{array}                                            
\usepackage{longtable}                                        
\usepackage{calc}                                             
\usepackage{multirow}                                         
\usepackage{hhline}                                          
\usepackage{ifthen} 
% \usepackage{babel}
\usepackage[T1]{fontenc}
\usepackage[most]{tcolorbox}
\usepackage{lmodern}
%\usepackage{tcolorbox}
\newtcolorbox{mybox}{colback=red!5!white,colframe=red!75!black}
%\newtcolorbox{mybox}{colback=grey!5!white,colframe=red!7https://www.overleaf.com/project/6215193db2c04bccf12b72395!black}
\usepackage{setspace}
\usepackage{verbatim}
\usepackage{etaremune}
\usepackage{url}
\usepackage{color}
\usepackage{graphicx}
\usepackage{natbib}
\usepackage{color}
\usepackage{array}
\usepackage{float}
\usepackage{wrapfig}
\usepackage[noend]{algpseudocode}
\usepackage{physics}
\usepackage{amssymb}
\usepackage{imakeidx}
\usepackage{hyperref}
\usepackage{lscape}
\makeindex

%\documentclass{standalone}%scrreprt
%\usepackage{pgfplots}
%\usepackage{tikz}
%\usepackage{tikz-network}
%\usepackage{amsmath}
%\usepackage{verbatim}
%\usepackage[dvipsnames]{xcolor}

%\usetikzlibrary{arrows,shapes,backgrounds,positioning}

\begin{comment}
Multilayer example following tikz-networks package
\end{comment}
%\usetikzlibrary{3d}

\begin{document}

\vspace{0.1 in}
\begin{equation}
\begin{pmatrix}
G_1(t) \\
G_2(t) \\
\vdots \\
G_l(t)
\end{pmatrix}
=
\begin{pmatrix}
Y_1(t) \\
Y_2(t) \\
\vdots \\
Y_l(t)
\end{pmatrix}
\circ
\begin{pmatrix}
L_1(t) \\
L_2(t) \\
\vdots \\
L_l(t)
\end{pmatrix}
\label{eq:6}
\end{equation}
\vspace{0.2 in}
\begin{equation}
\begin{pmatrix}
V_1(t) \\
V_2(t) \\
\vdots \\
V_l(t)
\end{pmatrix}
=
\begin{pmatrix}
E_{11}(t) & E_{12}(t) & \dots & E_{1l}(t) \\
E_{21}(t) & E_{22}(t) & \dots & E_{2l}(t) \\
\vdots & \vdots & \ddots & \vdots \\
E_{l1}(t) & E_{l2}(t) & \dots & E_{ll}(t)
\end{pmatrix}
\begin{pmatrix}
G_1(t) \\
G_2(t) \\
\vdots \\
G_l(t),
\end{pmatrix}
%\nonumber
\end{equation}
\vspace{0.2 in}
\hspace{12 in}\begin{equation}
\begin{pmatrix}
Z_1(t) \\
Z_2(t) \\
\vdots \\
Z_n(t)
\end{pmatrix}
=
\begin{pmatrix}
B_{11}(t) & B_{12}(t) & \dots & B_{1l}(t) \\
B_{21}(t) & B_{22}(t) & \dots & B_{2l}(t) \\
\vdots & \vdots & \ddots & \vdots \\
B_{n1}(t) & B_{n2}(t) & \dots & B_{nl}(t)
\end{pmatrix}
\begin{pmatrix}
V_1(t) \\
V_2(t) \\
\vdots \\
V_l(t),
\end{pmatrix}
%\nonumber
\end{equation}
\end{document}




